\documentclass[11pt]{article}
\usepackage{amsmath,amssymb,amsthm,booktabs,hyperref}
\newtheorem{theorem}{Theorem}
\newtheorem{lemma}[theorem]{Lemma}
\newtheorem{proposition}[theorem]{Proposition}
\title{Enumeration of bipartite graphs of degree five and diameter eight}
\author{human\_compression project}
\date{}
\begin{document}
\maketitle

\paragraph{Scope and attribution.}
The original problem is the census of connected graphs of order $32$,
degree $5$, girth $4$ and diameter $8$, considered by Cambie, Goedgebeur,
Jooken and Van den Eede in \emph{On the order-diameter ratio of
girth-diameter cages}, arXiv:2511.21144v3, Appendix~B, Table~6.
The present argument concerns its bipartite part. We use the separate
fourteen-skeleton theorem for this part: its source has SHA256
\texttt{e9bed0654c3aad3c6d3a90ee63a3dcca6c2cf8b1a9ce01f5afb168979b301ae4}.
Its proof and attribution are retained separately, rather than repeated
here. The following deductions and the proposed enumeration algorithm
require independent review. No ordinary class count, novelty claim,
solution of an open conjecture, or Lean formalization is asserted.

\paragraph{The normal form used below.}
Let $G$ be connected, simple, bipartite and $5$-regular on $32$ vertices,
and suppose that $a,b,c$ occur at positions $0,4,8$ on an eight-geodesic.
Write $A=N(a)$, $B=N(b)$ and $C=N(c)$. These are disjoint five-element
sets in one color class. Its remaining vertex is $z$. The other color
class consists of $a,c$, two sides $L,R$, and a core $X$ containing $b$.
Orient the endpoints so that all neighbors of $z$ outside $X$ lie in $L$.
The normal-form theorem states that $X$ consists of a class $X_0$ of
$t+1$ vertices with neighborhood $B$ and classes $X_i$, one for each
port $b_i\in B$, of size $m_i$ and neighborhood
$(B\setminus\{b_i\})\cup\{z\}$. Put $u=\sum_i m_i$ and $r=t+u$.
Then $u>0$, $r\leq3$, and $|L|+|R|=13-r$, with each side of size five
or six. The middle and port data have exactly the fourteen possibilities
in the retained table.

Each port meets only its assigned side, in $4-r+m_i$ vertices. There are
$5-u$ star incidences in $L$. If $h(v)$ counts the port incidences and
the possible star incidence of a side vertex $v$, then its five endpoint
rows omit $v$ exactly $h(v)$ times. For a side of size five these are
singleton omissions; for size six the omissions form a loopless
multigraph of five edges with degree sequence $h$. This is a necessary
and sufficient construction. Every permitted choice is connected, with
girth four and diameter eight.

\section*{Local enumeration}

Encode a side vertex by its incidence mask on its $b_D$ ports, together
with a star bit on $L$. There are at most four bits. A side table up to
permuting its vertices is precisely a multiset of $n_D\leq6$ masks,
with the prescribed column sums. If $c_e$ is the multiplicity of mask $e$
and $a_j$ is a column sum, its equations are
\[
 \sum_e c_e=n_D,\qquad
 \sum_e c_e\,\operatorname{bit}_j(e)=a_j,\qquad c_e\geq0.
\]
Every solution gives exactly such a multiset and every side table gives
a solution. Thus recursion over these nonnegative counts is exhaustive.
A dynamic program indexed by the next mask, the remaining number of
vertices and the remaining column sums has at most
\[
 (2^d+1)(n_D+1)\prod_j(a_j+1)
 \leq17\cdot7\cdot5^4=74\,375
\]
states. Its counts also provide exact ordinal intervals without
reconstructing earlier solutions.

On a side of size five no endpoint choice remains. On size six enumerate
the nondecreasing five-tuples of the fifteen possible pairs. Their number
is $\binom{19}{5}=11\,628$. Retain exactly the tuples with degree sequence
$h$. Every loopless multigraph of five edges appears once, with repeated
edges allowed and no ordering attached to the five endpoint rows.

\begin{lemma}[Local identification]
For a fixed side, order its ports by their missing multiplicities $m_i$.
Minimize over permutations of ports with equal multiplicity. For each
such ordering, sort side vertices by their incidence masks. For size six
also minimize its sorted omission-pair list over permutations within
equal-mask classes. The resulting least pair of mask list and omission
list is a complete invariant of this local incidence object, allowing
permutations of its side vertices, endpoint rows and ports of equal $m_i$.
\end{lemma}
\begin{proof}
A port permutation transforms the mask of every side vertex. Sorting the
new masks fixes the order of unequal masks. Any further permutation
preserving that list is exactly a permutation within equal-mask classes.
The unordered omission list already removes endpoint-row permutations.
This enumerates every remaining permitted action, so isomorphic local
objects have the same minimum. Conversely equality supplies explicit
port, side and row bijections preserving every incidence. On size five
the singleton omission multiset is determined by the mask population
counts, so its separate minimization has no effect.
\end{proof}

These local bijections extend through the whole middle. A permutation of
ports preserving $m_i$ permutes the corresponding equal-sized classes
$X_i$, fixes $X_0$ and $z$, and preserves all central adjacencies. The
two sides may therefore be reduced independently before joining them.
All pairs of their local representatives, across all fourteen rows,
cover every graph in the stated bipartite class. Ordinary duplicates
across different roots and rows still have to be identified.

\section*{Intrinsic choices of centers}

\begin{lemma}\label{diametral}
In the color class containing $a,b,c$, all pairs at distance eight run
between the two false-twin classes of $a$ and $c$, and every such pair
has distance eight.
\end{lemma}
\begin{proof}
The normal-form distance proof gives pairwise distance two between
distinct core vertices. Every vertex in the center color is either in
the core, within two of some core vertex, or a false twin of an endpoint.
Each endpoint false twin is at distance four from \emph{every} core
vertex. Consequently a pair other than two endpoint false twins has
distance at most six. Twins of the same endpoint have distance two.
Twins of opposite endpoints have the same distance eight as $a,c$.
\end{proof}

\begin{proposition}\label{roots}
The full quotient by equality of open neighborhoods has at most two
unordered pairs at distance eight. For each such pair, at most four
quotient vertices are at distance four from both endpoints. These give
all choices needed for ordinary, uncolored identification.
\end{proposition}
\begin{proof}
Lemma~\ref{diametral} gives one pair in the center color. If the other
color contains a pair at distance eight, choose an eight-geodesic between
them and apply the same normal-form theorem with these endpoints.
Lemma~\ref{diametral} then gives one pair in that color as well. Distances
between distinct twin classes equal distances between any of their
vertices: quotient paths lift because class adjacency is complete, and
graph paths project to quotient walks.

For a fixed endpoint pair every core vertex is at distance four from
both endpoints. Each side vertex is at distance two from its own
endpoint, while an endpoint has distance zero from itself. Opposite-color
distances are odd. Thus the simultaneous distance-four vertices are
exactly $X$. Its twin classes are $X_0$ and the nonempty $X_i$, numbering
at most $1+u\leq4$. Choosing another representative within a twin class
does not matter, since swapping twins is an automorphism fixing all
other vertices. This accounts for every possible center triple.
\end{proof}

\begin{theorem}[Ordinary canonical code]
For each endpoint-pair and midpoint-class choice of
Proposition~\ref{roots}, recover its normal form and orient its star side
as $L$. Form the tuple
\[
 (1,q,t,u,n_L,n_R,m_L,m_R,\mathcal L,\mathcal R),
\]
where each missing-multiplicity list is sorted, and $\mathcal L,\mathcal R$
are its local canonical codes. Let $C(G)$ be the least such tuple in a
fixed total order. Then two graphs in the stated class are ordinarily
isomorphic if and only if their codes $C(G)$ are equal.
\end{theorem}
\begin{proof}
An ordinary graph isomorphism preserves full twin classes, their weights,
distances, diametral endpoint pairs and simultaneous distance-four
classes. It therefore maps the entire set of eligible center choices
bijectively to that of the image graph. After choosing representatives,
its remaining changes are precisely the local symmetries already
minimized. It follows that the two sets of tuples, and hence their
minima, agree. This also covers an interchange of the two colors.

Conversely the tuple reconstructs the entire graph: expand $X_0$ and
$X_i$ with their specified neighborhoods, the two sides from their mask
lists, and the endpoint rows from their singleton or pair omissions.
Equality of tuples therefore gives a common graph to which both inputs
are isomorphic. The port and side permutations attaining the minimum,
the sorted omission rows and arbitrary bijections within twin classes
give explicit bijections on all thirty-two vertices. Composing those
bijections gives the required ordinary isomorphism.
\end{proof}

\paragraph{Finite bounds.}
The fourteen rows give $|H_L|+|H_R|\leq7$ for the two equal-multiplicity
port groups: the largest individual group permutes three equal ports.
A nonzero side mask occurs at most four times, since any of its set bits
has column sum at most four. Zero-mask vertices have omission degree
zero, so their permutations act trivially on the omission list. For a
side of size six the effective equal-mask group thus has order at most
$4!2!=48$. Together with the at most eight center choices, an elementary
canonical-code computation uses at most $8\cdot7\cdot48=2\,688$ small
local omission-list comparisons. This is an operation bound, not a
benchmark or proof that a particular resource-limited job will finish.

\paragraph{Certificates and remaining work.}
The proposed implementation emits finite side-profile intervals and
finite intervals of pairs of accepted local forms. A completed interval
must include every unit in its declared range; a bounded interruption
leaves the next unit explicit and uncredited. For every ordinary code it
also supplies a permutation from the constructed thirty-two-vertex graph
to the expansion of that code. Checking the entire adjacency relation
verifies each claimed identification independently.

Completeness still requires actual coverage of all side intervals and
all joined local pairs, independent review of the Source and lemmas,
and identification of the full ordinary union. A matching numerical
count or a rolling hash alone does not establish this. The existing
nonbipartite part of the original problem is separate. No current native
queue record is paid by this mathematical proposal, and no computation,
control, TeX compilation or Lean verification has been performed here.
\end{document}
